\chapter{The Journey of Code: From Human Text to Silicon Execution}

\section{The Four Magical Stations}

\begin{kidbox}{The Toy Factory Assembly Line}
Imagine your computer is a magical toy factory with four rooms. When you write a program on your screen, your text travels through each room:
\begin{enumerate}
    \item \textbf{Room 1: The Word Chomper (The Lexer)}: It takes your big sentence of letters and chomps it into little word-candy pieces called \textbf{Tokens}.
    \item \textbf{Room 2: The Master LEGO Builder (The Parser)}: It takes all those little candies and builds a beautiful, branching LEGO tree called the \textbf{Abstract Syntax Tree} (or \textbf{AST}).
    \item \textbf{Room 3: The Kind Teacher (The Linter)}: Before anyone runs around, the teacher inspects the LEGO tree to make sure you didn't connect pieces upside down or forget a wheel!
    \item \textbf{Room 4: The Engine Room (The Execution Engines)}: The factory workers look at the LEGO tree and make the gears turn, either by walking along the tree or translating it into secret robot cards!
\end{enumerate}
\end{kidbox}

\section{Station 1: The Lexer (Tokenizer)}

The Lexer scans the raw UTF-8 characters of your source file and categorizes them into discrete atomic units called \textbf{Tokens}.

\begin{lstlisting}[language=shaelang]
let score = 100 + 50;
\end{lstlisting}

When the Lexer processes this text, it emits a stream of tokens:
\begin{center}\vspace{8pt}
\centering
\begin{tabular}{lll}
\toprule
\textbf{Token Kind} & \textbf{Lexeme} & \textbf{Source Span (Line, Col)} \\
\midrule
\texttt{KeywordLet} & \texttt{"let"} & Line 1, Col 1 \\
\texttt{Identifier} & \texttt{"score"} & Line 1, Col 5 \\
\texttt{Assign} & \texttt{"="} & Line 1, Col 11 \\
\texttt{IntLiteral} & \texttt{"100"} & Line 1, Col 13 \\
\texttt{Plus} & \texttt{"+"} & Line 1, Col 17 \\
\texttt{IntLiteral} & \texttt{"50"} & Line 1, Col 19 \\
\texttt{Semicolon} & \texttt{";"} & Line 1, Col 21 \\
\bottomrule
\end{tabular}
\par\vspace{4pt}\small\textbf{Table: Token stream produced by the Shae Lexer}
\end{center}

\begin{deepbox}{Deterministic Finite Automata (DFA) and Spans}
In \texttt{src/lexer.rs} (and the self-hosting compiler \texttt{lib/compiler/lexer.shae}), tokenization is implemented as an explicit state machine over Unicode scalar values.
Numbers support arbitrary bases with prefix indicators:
\begin{itemize}
    \item Hexadecimal: \texttt{0xFF00} ($\text{radix}=16$)
    \item Binary: \texttt{0b1010\_0101} ($\text{radix}=2$ with underscore digit separators)
    \item Octal: \texttt{0o755} ($\text{radix}=8$)
\end{itemize}
Every token carries a cryptographic-like tracking structure called a \texttt{Span}:
\begin{lstlisting}[language=shaelang]
pub struct Span {
    pub start_line: Int,
    pub start_col: Int,
    pub end_line: Int,
    pub end_col: Int,
}
\end{lstlisting}
This allows Shae to produce beautiful compiler diagnostics with exact squiggly underlines.
\end{deepbox}

\section{Station 2: The Parser and Pratt Precedence}

\begin{kidbox}{Why Math Needs Rules: 2 + 3 * 4}
If you ask someone: "What is two plus three times four?", what is the answer?
Is it $(2 + 3) \times 4 = 20$? Or is it $2 + (3 \times 4) = 14$?
In math, multiplication is "stickier" than addition! It hugs its neighbors tighter.
The Parser knows all these secret sticking powers! When it builds the LEGO tree, it puts the times-four tightly next to three, and adds two at the very end.
\end{kidbox}

\begin{deepbox}{Top-Down Operator Precedence (Pratt Parsing)}
Shae uses a Pratt Parser (\texttt{src/parser.rs} and \texttt{lib/compiler/parser.shae}) to parse expressions with arbitrary operator precedence without suffering from the grammar bloat of pure recursive descent.

Each token has two binding powers:
\begin{enumerate}
    \item \textbf{Prefix Binding Power}: How tightly the operator binds when appearing before an expression (e.g., unary minus \texttt{-x}, logical not \texttt{!flag}, bitwise not \texttt{\~{}mask}).
    \item \textbf{Infix Binding Power}: How tightly an operator binds its left and right operands (e.g., \texttt{+}, \texttt{*}, \texttt{|>}, \texttt{==}).
\end{enumerate}

\begin{center}\vspace{8pt}
\centering
\begin{tabular}{lll}
\toprule
\textbf{Precedence Level} & \textbf{Operators} & \textbf{Associativity} \\
\midrule
1 (Lowest) & Assignment (\texttt{=}, \texttt{+=}, \texttt{-=}) & Right-to-Left \\
2 & Pipeline (\texttt{|>}) & Left-to-Right \\
3 & Logical OR (\texttt{or}, \texttt{||}) & Left-to-Right \\
4 & Logical AND (\texttt{and}, \texttt{\&\&}) & Left-to-Right \\
5 & Bitwise OR (\texttt{|}), XOR (\texttt{\^{}}) & Left-to-Right \\
6 & Bitwise AND (\texttt{\&}) & Left-to-Right \\
7 & Equality (\texttt{==}, \texttt{!=}) & Non-associative \\
8 & Relational (\texttt{<}, \texttt{<=}, \texttt{>}, \texttt{>=}) & Non-associative \\
9 & Shifts (\texttt{<<}, \texttt{>>}) & Left-to-Right \\
10 & Additive (\texttt{+}, \texttt{-}) & Left-to-Right \\
11 & Multiplicative (\texttt{*}, \texttt{/}, \texttt{\%}) & Left-to-Right \\
12 (Highest) & Unary (\texttt{-}, \texttt{!}, \texttt{\~{}}), Calls (\texttt{()}), Indexing (\texttt{[]}) & Right-to-Left \\
\bottomrule
\end{tabular}
\par\vspace{4pt}\small\textbf{Table: Shae Pratt Parser Precedence Table}
\end{center}
\end{deepbox}

\section{Station 3: The Static Linter and Semantic Analyzer}

Before code executes, the Shae Linter (\texttt{src/linter.rs}) traverses the AST to verify correctness:
\begin{itemize}
    \item \textbf{Unused Variable Tracking}: Detects bindings declared but never referenced. If a variable is intentionally unused, prefixing with an underscore (e.g., \texttt{\_unused}) silences the warning.
    \item \textbf{Dead Code Detection}: Flags statements that appear immediately after unconditional return statements, loops that never exit, or branches that cannot be reached.
    \item \textbf{Shadowing Diagnostics}: Alerts the developer when an inner scope variable accidentally shadows a variable in an outer scope or a built-in standard library symbol.
\end{itemize}

\section{Station 4: Dual-Engine Execution Parity}

One of Shae's most powerful architectural features is its \textbf{Dual-Engine Architecture}:
\begin{itemize}
    \item \textbf{The Tree-Walking Evaluator (\texttt{src/evaluator.rs})}: Directly walks the AST. Excellent for instant REPL execution, debugging, and interactive introspection.
    \item \textbf{The Stack-Based Bytecode VM (\texttt{src/vm.rs})}: Compiles AST down into a compact stream of \texttt{OpCode} bytes. Runs at high speed in a tight loop with linear instruction caching.
\end{itemize}
Every single test in Shae's test suite runs against both engines to guarantee 100\% behavioral parity!
